% Quarantined proof for the arbitrary-radius boundary-resonance output.

\begin{proof}[Proof of
\cref{prop:p062-nonstrict-boundary-resonance-counterexample}]
For \(a>0\), let \(\eta_a\) denote the \(a\)-Gauss point of the analytic
affine line.  The analytic closed immersion
\[
 (\mathbf G_{\mathrm m,K})^{\an}\longrightarrow(\mathbf A^2_K)^{\an},
             \qquad U\longmapsto(U,U^{-1}),
\]
sends \(\eta_R\) to a point \(y\in P\) at which \(TS-1=0\).  Hence
\(y\in W\), and \(W\) is nonempty.  It is relatively open by continuity
of the function \(|TS-1|\).

We next prove that every point of \(W\) lies over the \(R\)-Gauss point in
the first coordinate.  Let \(x\in W\).  The ultrametric inequality gives
\[
                         |T(x)S(x)|=1.
\]
Together with the two inequalities defining \(P\), this forces
\begin{equation}\label{eq:p062-resonant-coordinate-equalities}
                         |T(x)|=R,
                         \qquad |S(x)|=R^{-1}.
\end{equation}
If \(i\ne j\) and \(a_i,a_j\in K^\times\), then
\[
 |a_i|R^i=|a_j|R^j
       \quad\Longrightarrow\quad
 R^{i-j}\in|K^\times|,
\]
which is impossible because \(R\notin\sqrt{|K^\times|}\).  Thus the
numbers \(|a_i|R^i\), for the nonzero terms of a polynomial
\(f=\sum_i a_iT^i\), are pairwise distinct.  The largest term is unique,
so the strong triangle inequality and
\eqref{eq:p062-resonant-coordinate-equalities} give
\begin{equation}\label{eq:p062-resonant-gauss-restriction}
                         |f(x)|=\max_i|a_i|R^i.
\end{equation}
Therefore the image of \(x\) under the first-coordinate projection
\(
 \pi_T:(\mathbf A^2_K)^{\an}\to(\mathbf A^1_K)^{\an}
\)
is precisely \(\eta_R\).

The affine plane is boundaryless over \(K\).  Berkovich's relative-interior
composition formula and his identification of the relative interior of an
analytic domain with its topological interior
\cite[Proposition~1.5.5(i)--(ii)]{Berk93} therefore identify
\(\operatorname{Int}(P)\) with the ordinary interior of \(P\) in
\((\mathbf A^2_K)^{\an}\).  Suppose that \(x\in W\) belonged to this
interior.  There would be an ordinary open neighbourhood \(O\) of \(x\)
in the analytic affine plane with \(O\subset P\).  The projection
\(\pi_T\) is smooth and hence open by
\cite[Corollary~3.7.4]{Berk93}.  Consequently \(\pi_T(O)\) would be an
open neighbourhood of \(\eta_R\) contained in the closed disc
\(\{|T|\leq R\}\).  This is impossible: the map
\[
                    a\longmapsto\eta_a
\]
is continuous, as follows directly by evaluating each polynomial with its
Gauss norm, and \(\eta_a\to\eta_R\) through points with \(a>R\).  We have
proved
\begin{equation}\label{eq:p062-resonant-tube-in-boundary}
                    W\cap\operatorname{Int}(P)=\varnothing.
\end{equation}

The function in \eqref{eq:p062-resonant-cutoff} is continuous.  Moreover,
\[
 \operatorname{supp}(h_c)
       \subseteq\{|TS-1|\leq c\}
       \subset W,
\]
because \(c<1\).  It follows from
\eqref{eq:p062-resonant-tube-in-boundary} that \(h_c\) restricts to zero
on \(\operatorname{Int}(P)\), which proves
\eqref{eq:p062-resonant-zero-current}.  Finally,
\[
                         h_c(y)=1,
                         \qquad h_c(0,0)=0.
\]
Thus \(h_c\) is nonconstant, completing the counterexample.
\end{proof}

\begin{proof}[Proof of
\cref{prop:p062-nonstrict-scalar-stable-repair}]
Over \(L\), every defining radius belongs to \(|L^\times|\).  Hence
\(P_L\) is one of the strict endpoint polycylinders covered by
\cref{thm:p062-endpoint-polycylinder-v4-research}.  Applied to
\eqref{eq:p062-scalar-stable-current}, that theorem gives real constants
\(c,\lambda_1,\ldots,\lambda_q\) such that
\begin{equation}\label{eq:p062-scalar-stable-upstairs-formula}
          h_L=c+\sum_{j=1}^q\lambda_j\log|S_j|
          \qquad\text{on }P_L.
\end{equation}

The scalar-extension map \(p_L:P_L\to P\) is surjective
\cite[\S11.1.9]{CLD25}.  Given \(x\in P\), choose
\(x_L\in p_L^{-1}(x)\).  Scalar extension preserves the seminorms of the
coordinate functions, so
\[
 \begin{aligned}
 h(x)
   &=h_L(x_L)\\
   &=c+\sum_{j=1}^q\lambda_j\log|S_j(x_L)|\\
   &=c+\sum_{j=1}^q\lambda_j\log|S_j(x)|.
 \end{aligned}
\]
This proves \eqref{eq:p062-scalar-stable-affine}.  Uniqueness follows by
restriction to the standard Gauss skeleton of the nonempty open
polyannulus core.  The functions
\[
                  1,\log|S_1|,\ldots,\log|S_q|
\]
are affinely independent on its open logarithmic box.  When \(q=0\),
evaluation at any point gives uniqueness of \(c\).

For completeness, a strictifying extension always exists.  Starting from
\(K\), adjoin one indeterminate for each of the finitely many defining
radii, endow the successive rational-function fields with the corresponding
Gauss absolute values, and complete after the last step.  The resulting
complete extension has an element of each prescribed absolute value.  The
substantive extra hypothesis is therefore not the existence of \(L\), but
the upstairs current equation
\eqref{eq:p062-scalar-stable-current}.
\end{proof}
